\chapter*{Введение}
\addcontentsline{toc}{chapter}{Введение}

Современные компьютерные сети представляют собой сложные распределённые системы, обеспечивающие передачу данных между миллионами устройств по всему миру. С ростом объёмов передаваемой информации и увеличением числа подключённых устройств возрастает и количество угроз информационной безопасности. Ежегодно регистрируются тысячи инцидентов, связанных с DDoS-атаками, сканированием портов, распространением вредоносного программного обеспечения, попытками несанкционированного доступа и утечками конфиденциальных данных.

Традиционные системы обнаружения вторжений (IDS), такие как Snort, Suricata и Zeek, используют сигнатурный метод анализа, который эффективен против известных атак, но неспособен выявлять новые, ранее не встречавшиеся угрозы (zero-day attacks). Это обуславливает необходимость разработки систем, способных обнаруживать аномалии на основе анализа поведенческих характеристик сетевого трафика с применением методов машинного обучения.

Данный проект посвящён разработке платформы выявления аномалий сетевой активности на основе потокового анализа сетевых данных. Платформа реализует полный цикл обработки: от захвата сетевого трафика до детектирования аномалий и визуализации результатов.

\textbf{Цель работы:} разработать программную платформу для обнаружения аномалий сетевой активности, использующую методы машинного обучения и потоковой обработки данных.

\textbf{Задачи работы:}
\begin{enumerate}
    \item Провести анализ существующих подходов к мониторингу и выявлению аномалий сетевой активности, сформулировать требования к разрабатываемой платформе.
    \item Спроектировать архитектуру системы потоковой обработки сетевых данных.
    \item Реализовать модуль сбора и потоковой обработки сетевых событий.
    \item Разработать и интегрировать алгоритмы выявления аномалий сетевой активности.
    \item Провести тестирование системы на синтетических сетевых данных и проанализировать полученные результаты.
\end{enumerate}

\textbf{Практическая значимость} работы заключается в создании инструмента, позволяющего автоматизировать процесс обнаружения сетевых аномалий, что может быть использовано как в учебных целях, так и в качестве основы для корпоративных систем мониторинга безопасности.

Отчёт состоит из 6 глав, содержит 12 рисунков, 8 таблиц и 42 страницы.
